\section{Tessellations of the sphere}\label{tessellations-of-the-sphere}

To begin our journey into non-euclidean geometries we will first look at
geometry of the sphere. Understanding spherical geometry will give us
helpful insights to understand and interpret hyperbolic geometry.

\subsection{Geometry on the sphere}\label{geometry-on-the-sphere}

\begin{itemize}
\tightlist
\item[$\square$]
  Geodesics (what is the shortest path on the sphere)
\item[$\square$]
  Antipodal
\item[$\square$]
  Dihedral Angle
\end{itemize}

\subsubsection{Great Circles}\label{great-circles}

\begin{itemize}
\tightlist
\item[$\square$]
  Todo, definition for the unit sphere \(S\) and representation of
  planes by unit normal vectors \(\hat n\)
\end{itemize}

A great circle is given by the intersection of the sphere with a plane
through it's center. For any two points on the great circle the arc
formed between them is the shortest path the sphere.

\begin{itemize}
\tightlist
\item[$\square$]
  Show that this intersection is actually a circle?
\end{itemize}

Let \(A, B\) be two points on the sphere. The normal vector of a great
circle through \(A\) and \(B\) is given by the cross product
\(\vec{A} \times \vec{B}\).

\begin{itemize}
\tightlist
\item[$\square$]
  We proof that arcs of great circles form the shortest path between two
  points
\item[$\square$]
  Note that normal vectors on the xy-plane form great circles through
  the north and south pole of the sphere.
\item[$\square$]
  The perpendicular line to a great circle through it's center is called
  it's axis
\item[$\square$]
  The two point intersection of a great circles axis with the sphere are
  called it's poles.
\item[$\square$]
  A great circle splits the sphere into two hemispheres
\end{itemize}

Let \(x \in S\) be the elements of a sphere and \(\hat{n}\) the unit
normal vector of a great circle. The great circle partitions the sphere
into two hemispheres:

\begin{itemize}
\tightlist
\item
  A positive hemisphere \(S_+ = \set{x \in S: x \cdot \hat{n} > 0}\),
\item
  a negative hemisphere \(S_- = \set{x \in S: x \cdot \hat n < 0}\),
\end{itemize}

\subsubsection{Spherical Lunes}\label{spherical-lunes}

\begin{itemize}
\tightlist
\item[$\square$]
  Special case of a digon
\item[$\square$]
  Polygon with two vertices (possible in S, degenerate in E)
\end{itemize}

Let \(g_1, g_2\) be two great circles and \(n_1, n_2\) their respective
unit normal vectors. The two cricles intersect at an angle \(\theta\)
given by \(cos(\theta) = \hat n_1 \cdot \hat n_2\). Let
\(S_{1+}, S_{2+}\) be the positive hemisphere given by the great
circles. We define a spherical lune \(L\) to be the wedge shaped
intersection of the two hemispheres \(L = S_{1+} \cap S_{2+}\).

The area of a lune is given by it's angle \(A(L_\theta) = 2\theta r^2\).
A proof is left as an exercise to the reader :) This fact will proof
useful later.

Note that the intersection \(L’ = S_{1-} \cap S_{2-}\) is a congruent
antipodal lune to \(L\). Meaning \(A(L_\theta) = A(L_\theta’)\).

\begin{itemize}
\tightlist
\item[$\square$]
  Not sure if we should add this. It is rigorous but also obfuscating
  the point a lot

  \begin{itemize}
  \tightlist
  \item[$\square$]
    \(\overline{L} = \overline{S_{1+}} \cap \overline{S_{2+}}\) lune
    including boundary
  \item[$\square$]
    \(A(L) = A(\overline{L})\) including boundary does not change area
  \end{itemize}
\end{itemize}

\paragraph{Spherical Triangles}\label{spherical-triangles}

Let \(A, B, C \in S\) be points on the unit sphere forming the vertices
of a spherical triangle. It's edges \(a,b,c\) are given by great circles
with respective unit normal vectors \(\hat n_a, \hat n_b, \hat n_c\)
with

\[\begin{split}
n_a = C \times B \\
n_b = A \times C \\
n_c = B \times A
\end{split}
\]

It is good to be rigorous at this stage with how we define the great
circles that form our triangle. We want
\(\triangle = S_{a+} \cap S_{b+} \cap S_{c+}\). This will proof useful
once we get into algorithms to generate tilings. Note this also has to
be considered in Euclidean geometry when we represent lines in their
normal form.

\[\begin{split}
\alpha= n_b \cdot n_c \\
\beta = n_c \cdot n_b \\
\gamma = n_a \cdot n_b
\end{split}\]

The angle sum of a spherical triangle is greater than \(\pi\). There is
an elegant proof to show this.

The unit normal vectors are non-colinear and non-planar if and only if
A, B, C span a triangle. This is exactly the case when the scalar triple
product \([n_a, n_b, n_c] = n_a \cdot (n_b \times n_c) > 0\).

Let us construct the spherical lunes \(L_\alpha, L_\beta, L_\gamma\) and
their antipodal counterparts \(L_\alpha’, L_\beta’, L_\gamma’\). Note
that each lune contains either the original triangle \(\triangle\) or
it's antipodal counterpart \(\triangle’\). Note further that we can
reconstruct the sphere by forming the unions of all lunes
\(\bigcup \overline{L} = S\).

Summing up the area of all lunes we get back the area of the sphere,
however overcounting by four areas of our original triangle. Note our
lemma \(A(L)=A(L’)\) hence \(A(L)+A(L’) = 2A(L)\). Simplifying we get an
equation for the angle sum.

\[\begin{split}
2A(L_\alpha) + 2A(L_\beta) + 2A(L_\gamma) &= 4 \pi r^2 + 4A(\triangle) \\
&\Downarrow \\
A(L_\alpha) + A(L_\beta) + A(L_\gamma) &= 2 \pi r^2 + 2A(\triangle) \\
&\Downarrow \\
2\alpha r^2 + 2\beta r^2 + 2\gamma r^2 &= 2\pi r^2 + 2A(\triangle) \\
&\Downarrow \\
\alpha + \beta + \gamma &= \pi + \frac{A(\triangle)}{r^2}
\end{split}\]

This gives us a nice equation for the angle sum in relation to the
triangles area. Let us consider the area bounds of a spherical triangle.
It is clear that the area of a triangle must be greater than 0 and
cannot exceed \(2\pi\). Therefore \(0 < T < 2\pi\)

Further we get a bound on the angle sum
\(\pi < \alpha + \beta + \gamma < 3\pi\).

This also gives us a formula for the area of a spherical triangle where

\[A(\triangle) = \alpha + \beta + \gamma - \pi\]

\subsection{Constructing regular
tesselations}\label{constructing-regular-tesselations}

\subsubsection{Hosohedron}\label{hosohedron}

\subsubsection{Triangle Tilings}\label{triangle-tilings}

\begin{itemize}
\tightlist
\item[$\square$]
  split hosohedron along equator
\item[$\square$]
  combinations for pqr with angle sum constraint
\end{itemize}

\subsubsection{Regular and Uniform
Tesselations}\label{regular-and-uniform-tesselations}

\begin{itemize}
\tightlist
\item[$\square$]
  Dual graphs on triangle tilings

  \begin{itemize}
  \tightlist
  \item[$\square$]
    already getting into kaleidoscope construction through symmetry
    groups?
  \end{itemize}
\end{itemize}

\subsection{Excursion Polyhedra}\label{excursion-polyhedra}

\begin{itemize}
\tightlist
\item[$\square$]
  Tilings of the sphere give us graphs of 3d polyhedra sharing same
  symmetries and duality

  \begin{itemize}
  \tightlist
  \item[$\square$]
    Platonic solids
  \item[$\square$]
    Prisms
  \end{itemize}
\end{itemize}
